% Standalone synthesis of the CHL, Nikulin, and Gritsenko-Cl\'ery rows.
% Source anchors:
%   Gritsenko-Nikulin, Automorphic forms and Lorentzian Kac-Moody
%     algebras II, Part II, Thm. 2.1.
%   Borcherds, Invent. Math. 132 (1998), Thm. 13.3.
%   Gritsenko, Arithmetical lifting and its applications, 1999.
%   Gritsenko-Cl\'ery, arXiv:0812.3962, Thms. 1.4 and 3.1.
%   Nikulin, finite symplectic automorphisms of K3 surfaces.
%   Oberdieck-Pandharipande for the K3 x E inverse-square scalar
%     normalisation; Bryan-Oberdieck for CHL quotient curve-counting.

\providecommand{\kBKM}{\kappa_{\mathrm{BKM}}}
\providecommand{\kch}{\kappa_{\mathrm{ch}}}
\providecommand{\kcat}{\kappa_{\mathrm{cat}}}
\providecommand{\kfib}{\kappa_{\mathrm{fiber}}}
\providecommand{\Z}{\mathbb{Z}}
\providecommand{\C}{\mathbb{C}}
\providecommand{\Sp}{\mathrm{Sp}}
\providecommand{\wt}{\mathrm{wt}}
\providecommand{\CoHA}{\mathrm{CoHA}}
\providecommand{\CY}{\mathrm{CY}}
\providecommand{\CHL}{\mathrm{CHL}}
\providecommand{\K}{K3}
\providecommand{\cO}{\mathcal{O}}
\providecommand{\Dfive}{D_5}
\providecommand{\Delfnm}[2]{\Delta_{#1}^{(#2)}}
\providecommand{\Yplus}{Y^{+}}
\providecommand{\Winf}{\mathcal{W}_{1+\infty}}
\providecommand{\gDelta}{\mathfrak{g}_{\Delta_5}}
\providecommand{\gBPS}{\mathfrak{g}_{\mathrm{BPS}}}
\providecommand{\Hall}{\mathrm{Hall}}

\section*{Eight Borcherds Rows and Their Recognition Data}
\label{note:eight-borcherds-rows-recognition-data}

\subsection*{Three Catalogues}

Three row systems meet at the primitive level-one form $\Delta_5$.
Away from that row their indices, input forms, and geometric hosts are
different.  The primitive compact CHL ladder is indexed by
$N\in\{1,2,3,4,6\}$; the non-CHL Nikulin boundary targets are indexed
by $N\in\{5,7,8\}$; the Gritsenko-Cl\'ery diagonal-divisor catalogue
is indexed by triples $(t,N;k)$.

\[
\begin{array}{c|c|c|c|l}
N & \text{system} & c_N(0) & \kBKM & \text{primitive denominator}\\\hline
1 & \CHL & 10 & 5 & \Delta_5\\
2 & \CHL & 8  & 4 & \Delfnm{4}{2}\\
3 & \CHL & 6  & 3 & \Delfnm{3}{3}\\
4 & \CHL & 4  & 2 & \Delfnm{2}{4}\\
6 & \CHL & 2  & 1 & \Delfnm{1}{6}
\end{array}
\]

This is the primitive CHL denominator ladder.  It is the
Gritsenko-Nikulin additive-lift tower at the compact CHL orders
$N\in\{1,2,3,4,6\}$.  Borcherds's singular-theta product theorem gives
the weight of the automorphic product as one half of the constant term
of the input weak Jacobi or vector-valued modular form; in this
normalisation,
\[
  \kBKM(\Phi_N)=\wt(\Phi_N)=\frac{c_N(0)}{2}.
\]
The physics Igusa convention squares the primitive denominator; hence
the physical weights are $(10,8,6,4,2)$, while the primitive BKM
weights are $(5,4,3,2,1)$.

\[
\begin{array}{c|c|c|l}
N & c_N(0) & \wt(F_N) & \text{geometric datum required}\\\hline
5 & 4 & 2 & \text{Borcea-Voisin CY}_3\ \text{row map}\\
7 & 2 & 1 & \text{singular or twisted CY}_3\ \text{descent}\\
8 & 2 & 1 & \text{Kummer-type CY}_4/\text{resolved CY}_3\ \text{bridge}
\end{array}
\]

These are non-CHL Nikulin boundary rows.  The displayed numbers are
automorphic target constants and weights:
$c_5(0)=4$, $c_7(0)=2$, $c_8(0)=2$.  They become
$\kBKM=(2,1,1)$ precisely when a denominator-recognition theorem supplies
a Hall source, a row map, and a common Borcherds denominator.  Order
$5$ requires a Borcea-Voisin comparison with reduced DT and Hall
data; order $7$ requires a singular or gerby descent theory; order $8$
requires a bridge between the holomorphic-symplectic fourfold theory
and a resolved CY$_3$ theory.

\[
\begin{array}{c|c|c|c|l}
(t,N;k) & F_{t,N} & \wt(F_{t,N}) & 2\wt(F_{t,N})
& \text{construction}\\\hline
(1,1;5) & \Delta_5 & 5 & 10 &
  \text{additive lift on }\Gamma_1\\
(2,1;2) & \Delta_2 & 2 & 4 &
  (1,2)\text{-polarised abelian surface}\\
(1,2;3) & \nabla_3 & 3 & 6 &
  \Gamma_0^{(2)}(2)\text{ Borcherds product}\\
(3,1;1) & \Delta_1 & 1 & 2 &
  (1,3)\text{-polarised abelian surface}\\
(1,3;2) & \nabla_2 & 2 & 4 &
  \Gamma_0^{(2)}(3)\text{ Borcherds product}\\
(4,1;\tfrac12) & \Delta_{1/2} & \tfrac12 & 1 &
  \text{genus-two theta-nullwert}\\
(1,4;\tfrac32) & \nabla_{3/2} & \tfrac32 & 3 &
  \Gamma_0^{(2)}(4)\text{ square-root product}\\
(2,2;1) & Q_1 & 1 & 2 &
  \Gamma_2(2)\text{ Borcherds product}
\end{array}
\]

This is the Gritsenko-Cl\'ery diagonal-divisor catalogue.  The triple
$(t,N;k)$ and the paramodular group $\Gamma_t(N)$ are part of the row
datum.  The row-wise weight formula is the Gritsenko-Cl\'ery
cusp-weighted Borcherds product formula,
\[
 \wt(F_{t,N})=
 \frac12\sum_{f/e\in\mathcal P(N)}
 \frac{h_e}{N_e}\,c_{f/e}(0,0).
\]
The row $(t,N;k)=(1,1;5)$ alone is simultaneously a CHL row.  In triple
order the weights are
\[
 (5,2,3,1,2,\tfrac12,\tfrac32,1),
\]
and the corresponding twice-weights are
\[
 (10,4,6,2,4,1,3,2).
\]

\subsection*{Level-One Normalisation}

The primitive K3-Igusa normalisation is
\[
 \Dfive = 64^{-1}\Delta_5,\qquad
 \Phi_{10}^{\mathrm{un}}=\Delta_5^2,\qquad
 \Phi_{10}^{\mathrm{OP}}=\Dfive^2=4096^{-1}\Delta_5^2 .
\]
The Oberdieck-Pandharipande inverse-square convention is
\[
 Z_{\mathrm{OP}}
 =-\bigl(\Phi_{10}^{\mathrm{OP}}\bigr)^{-1}
 =-D_5^{-2}
 =-4096\,\Delta_5^{-2}.
\]
The primitive BKM denominator is controlled by $\Delta_5$ and has
$\kBKM(\gDelta)=5$.  The Oberdieck-Pandharipande reduced-DT scalar
square is controlled by $\Phi_{10}^{\mathrm{OP}}$ and has weight $10$
before the normalising scalar is removed.  Scalar inverse-square
normalisation, primitive denominator recognition, and physical
partition-function interpretation are separate readings of the same
level-one product.

\subsection*{The K3 Times E Invariants}

For $Y=\K\times E$ the invariant split is
\[
\begin{array}{c|c|l}
\text{invariant} & \text{value} & \text{construction}\\\hline
\kcat(Y) & 0 &
  \chi(\cO_{\K})\chi(\cO_E)=2\cdot 0\\
\kch(Y) & 0 &
  \sum_q(-1)^q h^{0,q}(Y)=1-1+1-1\\
\kappa_{\mathrm{ch}}^{\mathrm{Heis}}(Y) & 3 &
  \text{Mukai-Heisenberg specialisation on the elliptic curve}\\
\kBKM(\gDelta) & 5 &
  c_1(0)/2=10/2\\
\kfib(Y) & 24 &
  \operatorname{rk}\widetilde H(\K,\Z)=24
\end{array}
\]
The values come from separate constructions.  The compact Hodge
supertrace, the Heisenberg specialisation, the Borcherds denominator,
and the Mukai-lattice fibre rank have distinct sources and distinct
functorial meanings.  Each value retains its own domain of definition:
$\kBKM(\gDelta)$ is a Borcherds weight, $\kch(Y)$ is the compact Hodge
supertrace, and $\kfib(Y)$ is a Mukai-lattice rank.

\subsection*{Hall-Drinfeld and Mukai Branches}

The BKM branch over $\K\times E$ is the Hall-Drinfeld branch.  Its
target is a completed double of a positive Hall half,
\[
  \mathcal D_\hbar\bigl(\Yplus_{\Hall}(\K\times E)\bigr)
  \leadsto \gDelta ,
\]
after orientation, PBW, and primitive-recognition data are supplied.
The Mukai Yangian branch is the self-mirror K3/Heisenberg branch.  It
uses Mukai currents and rational Yangian notation on the abelian
Heisenberg layer.  The two branches share K3 lattice numerics; their
generators, brackets, and recognition problems are different.

For toric $\C^3$ the CoHA is the positive half
$\CoHA(\C^3)=\Yplus(\widehat{\mathfrak{gl}}_1)$.  The corresponding
full $\Winf$ object belongs to the Drinfeld double, centre,
completion, or evaluation layer.  For a CY$_3$ output algebra $A$ of
$\Phi$, the native structure is $E_1$.  The $E_2$ structure appears
after passing to $Z(\mathrm{Rep}(A))$, to a double, or to a
module/centre layer constructed from $A$.

\subsection*{Recognition Data}

\begin{enumerate}
\item \emph{Primitive CHL rows.}  The automorphic constants and weights
are theorem-level data.  A compact Hall-Drinfeld/BKM identification
also needs orientation data on the reduced moduli, PBW comparison for
the BPS Lie algebra, and a primitive-recognition certificate matching
the Gritsenko-Nikulin bracket.
\item \emph{Order $5$.}  A Borcea-Voisin host must provide a reduced
DT theory, a Hall source, and a row map from that source to the
weight-$2$ Borcherds product.
\item \emph{Order $7$.}  The fixed-point quotient forces a singular or
twisted target.  The required data are a differential-character or
gerbe descent package, an orientation line on the twisted moduli, and
an index formula producing the weight-$1$ product.
\item \emph{Order $8$.}  The Mongardi-Tari-Wandel Kummer-type object
is naturally holomorphic-symplectic of dimension four.  The open datum
is the enumerative bridge: DT$_4$ on the fourfold, or a resolved
CY$_3$ theory with the same Borcherds product.
\item \emph{Gritsenko-Cl\'ery rows.}  The eight diagonal-divisor forms
are automorphic theorem-level objects.  A CY host claim for any row
other than $\Delta_5$ requires an explicit category, stability
condition, reduced virtual class, Hall product, and denominator
comparison.
\item \emph{Physical readings.}  The inverse-square partition function
uses the OP scalar normalisation above.  A black-hole or topological
string interpretation is a physical heuristic unless a reduced-DT
theorem, orientation datum, and duality frame are specified.
\end{enumerate}

\subsection*{Compact Synthesis}

The primitive CHL ladder has constants
$(10,8,6,4,2)$ and BKM weights $(5,4,3,2,1)$ at
$N\in\{1,2,3,4,6\}$.  The non-CHL Nikulin boundary targets have
automorphic constants $(4,2,2)$ and target weights $(2,1,1)$ at
$N\in\{5,7,8\}$; those target weights become $\kBKM$ precisely when
denominator recognition is supplied.  The Gritsenko-Cl\'ery
diagonal-divisor catalogue has weights
$(5,2,3,1,2,\tfrac12,\tfrac32,1)$ indexed by triples $(t,N;k)$, with
twice-weights $(10,4,6,2,4,1,3,2)$.  The level-one equalities
$\Dfive=64^{-1}\Delta_5$,
$\Phi_{10}^{\mathrm{OP}}=4096^{-1}\Delta_5^2$, and
$Z_{\mathrm{OP}}=-4096\Delta_5^{-2}$ anchor the common row.  All other
comparisons require row maps, normalisation choices, and host data
stated before any geometric or Hall-theoretic identification is
asserted.
